\uuid{yEfx}
\exo7id{470}
\titre{exo7 470}
\auteur{cousquer}
\organisation{exo7}
\datecreate{2003-10-01}
\isIndication{false}
\isCorrection{false}
\chapitre{Propriétés de R}
\sousChapitre{Maximum, minimum, borne supérieure}
\cpge{mpsi-mp2i.raisonnement-vocabulaire-ensembliste.a, pcsi.raisonnement-vocabulaire-ensembliste.a, ptsi.raisonnement-vocabulaire-ensembliste.a, mpsi-mp2i.calcul-algebrique-trigonometrie.c, pcsi.calcul-algebrique-trigonometrie.c, ptsi.calcul-algebrique-trigonometrie.c, bcpst1.logique-ensembles-raisonnement.a, bcpst1.nombres-reels}
\module{Analyse}
\niveau{L1}
\difficulte{}

\contenu{
\texte{
Étant donné un ensemble $A \subset\mathbb{R}$,
écrire avec des quantificateurs
les propriétés suivantes :
}
\begin{enumerate}
    \item \question{$10$ est un majorant de $A$,}
    \item \question{$m$ est un minorant de $A$,}
    \item \question{$P$ n'est pas un majorant de $A$,}
    \item \question{$A$ est majoré,}
    \item \question{$A$ n'est pas minoré,}
    \item \question{$A$ est borné,}
    \item \question{$A$ n'est pas borné.}
\end{enumerate}
}
